% Complete macroscopic marked equilibrium provider, 2026-09-14.
% Input fragment; its endpoint proof is endpoint_existence.tex in this folder.
\section{The actual macroscopic marked Gamma equilibrium}
\label{sec:macroscopic-marked-equilibrium}

The parameter in this section is $\lambda\geq0$, and $\kappa=2\lambda$.
The original packet centre remains $c=k/2$; the auxiliary cutoff
$\kappa$ is obtained by the explicit parameter maps
$K\mapsto L=(K-1)/4$, $(L,q)\mapsto L/q$, and
$\lambda\mapsto\kappa=2\lambda$.
The parameter $\lambda$ is the limit of the actual ratio $L/q$ in the marked factor with
$K=4L+1$; neither $K$ nor $L$ is identified with the original packet
parameter $k$.  Define, on the original limiting coordinate $x>0$,
\begin{align}\tag{KME1}
 h_\lambda^{\rm mark}(x)
 &=\int_0^\lambda\log(x+4t^2)\,dt\\
 &=\lambda\log(x+4\lambda^2)-2\lambda
          +\sqrt x\arctan(2\lambda/\sqrt x),\notag\\
 V_\lambda(x)&=\pi\sqrt x-2\log x-2h_\lambda^{\rm mark}(x).
 \notag
\end{align}
The equality is verified by differentiating the displayed antiderivative
with respect to $\lambda$ and evaluating its zero value at $\lambda=0$.
Every additive constant in this potential is retained.  In terms of $\kappa$,
\begin{equation}\tag{KME2}
 V_\lambda(x)
 =2\sqrt x\arctan(\sqrt x/\kappa)-2\log x
                  -\kappa\log(x+\kappa^2)+2\kappa,
 \qquad
 V'_\lambda(x)=\int_\kappa^\infty\frac{du}{x+u^2}-\frac2x.
\end{equation}
At $\kappa=0$ use $\arctan(\sqrt x/0)=\pi/2$.
For $\kappa>0$ the derivative is
$\arctan(\sqrt x/\kappa)/\sqrt x-2/x$; integration of
$1/(x+u^2)$ proves the same formula at $\kappa=0$ by its limiting value.

For a probability measure on $(0,\infty)$ put
\begin{equation}\tag{KME3}
 \mathcal E_\lambda(\mu)=\int V_\lambda\,d\mu
        -\iint\log|x-y|\,d\mu(x)d\mu(y),\qquad
 F_\lambda=-\inf_\mu\mathcal E_\lambda(\mu).
\end{equation}
This is the supremum convention for $F_\lambda$ used in the determinant
calculation. The potential tends to $+\infty$ at both ends, and
$V_\lambda(x)=\pi\sqrt x-(2+\kappa)\log x+O(x^{-1})$ as $x\to\infty$.
For $\kappa>0$, the integral of $\log(x+u^2)$ over $0<u<\kappa$ has a finite
limit as $x\downarrow0$; for $\kappa=0$ it vanishes identically. Thus the
coefficient of the divergent $-\log x$ term at zero is always exactly $2$.
The extended pair kernel
$\tfrac12V_\lambda(x)+\tfrac12V_\lambda(y)-\log|x-y|$
is bounded below, since it is at least $g(x)+g(y)$ with
$g(x)=V_\lambda(x)/2-\log(1+x)$ bounded below.
Finite energy forces $\int(\sqrt x+|\log x|)\,d\mu<\infty$ by the
same lower bound.  This specifies the energy on all probability measures,
with value $+\infty$ allowed, without subtracting divergent terms.

\input{endpoint_existence.tex}

\subsection{The positive density and the full variational inequality}
Use the unique endpoints $a=a(\kappa)$ and $b=b(\kappa)$ established in
(KEP1)--(KEP20), so $0<a<b$. Write
\begin{equation}\tag{KME4}
 m_{\rm end}=\frac{a+b}{2},\qquad d_{\rm end}=\frac{b-a}{2},\qquad
 D_s=\sqrt{(a+s)(b+s)},\quad R_u=D_{u^2},\quad D_0=\sqrt{ab}.
\end{equation}
The two actual endpoint equations are
\begin{equation}\tag{KME5}
 \int_\kappa^\infty\frac{du}{R_u}=\frac2{D_0},\qquad
 \int_\kappa^\infty\left(1-\frac{u^2}{R_u}\right)du=4.
\end{equation}
Define, for every $x>0$,
\begin{align}\tag{KME6}
 H_\kappa(x)
 &=\frac2{xD_0}-\int_\kappa^\infty\frac{du}{(x+u^2)R_u}\\
 &=\frac1x\int_\kappa^\infty
                  \frac{u^2\,du}{(x+u^2)R_u}>0,\notag\\
 \rho_\kappa(x)&=\frac{\sqrt{(b-x)(x-a)}}{2\pi}H_\kappa(x)
                  \quad(a<x<b),\notag
\end{align}
with density zero outside $[a,b]$.  The second identity is precisely the
first endpoint equation in (KME5); it proves strict positivity rather
than assuming it. All integrals converge absolutely. On a compact
positive $x$ interval their derivatives in $x$ converge locally uniformly:
the last integrand in (KME6) is $O(u^2)$ at zero and $O(u^{-2})$
at infinity. Thus $H_\kappa$ is smooth and strictly positive on $[a,b]$.

Let $\mathcal R(z)=\sqrt{(z-a)(z-b)}$ on
$\mathbb C\setminus[a,b]$, with $\mathcal R(z)/z\to1$ at infinity.
In particular $\mathcal R(x)$ is negative for $x<a$ and positive for
$x>b$, and its upper boundary value in the interval is
$i\sqrt{(b-x)(x-a)}$.
The elementary arcsine and semicircle integrals give
\begin{align}\tag{KME7}
 \frac1\pi\int_a^b
    \frac{dt}{(z-t)\sqrt{(b-t)(t-a)}}&=\frac1{\mathcal R(z)},\\
 \frac1\pi\int_a^b
    \frac{\sqrt{(b-t)(t-a)}}{z-t}\,dt&=z-m_{\rm end}-\mathcal R(z),\notag\\
 \frac1\pi\int_a^b
    \frac{\sqrt{(b-t)(t-a)}}{t+s}\,dt&=m_{\rm end}+s-D_s.\notag
\end{align}
To verify the first equation directly for $z>b$, put
$t=m_{\rm end}+d_{\rm end}\cos\theta$ and $w=\tan(\theta/2)$; its integral becomes
$\pi^{-1}\int_0^\infty2\,dw/[(z-b)+(z-a)w^2]$.
Analytic continuation proves the stated domain. Polynomial division
of $(b-t)(t-a)$ then proves the other two identities.
Partial fractions yield
\begin{equation}\tag{KME8}
 A_s(z):=\frac1\pi\int_a^b
  \frac{\sqrt{(b-t)(t-a)}}{(t+s)(z-t)}\,dt
   =1-\frac{D_s+\mathcal R(z)}{z+s}.
\end{equation}
At $z=-s$ the right side has its removable value.

The mass of the density is exactly
\begin{align}\tag{KME9}
 \int_a^b\rho_\kappa(t)\,dt
 &=\frac{m_{\rm end}-D_0}{D_0}
   -\frac12\int_\kappa^\infty\left(\frac{m_{\rm end}+u^2}{R_u}-1\right)du\\
 &=-1+\frac12\int_\kappa^\infty
                  \left(1-\frac{u^2}{R_u}\right)du=1.\notag
\end{align}
The cancellation of the terms containing $m_{\rm end}$ uses the first equation
in (KME5), and the last equality uses its second equation.
If $G_\kappa(z)=\int_a^b\rho_\kappa(t)/(z-t)\,dt$, the same substitutions give,
initially off the negative real axis and off $[a,b]$,
\begin{align}\tag{KME10}
 G_\kappa(z)&=\frac1{D_0}A_0(z)
                   -\frac12\int_\kappa^\infty\frac{A_{u^2}(z)}{R_u}\,du\\
 &=\frac12\left\{\int_\kappa^\infty\frac{du}{z+u^2}-\frac2z
                            -\mathcal R(z)H_\kappa(z)\right\}
   =\frac{V'_\lambda(z)-\mathcal R(z)H_\kappa(z)}2.\notag
\end{align}
All the exchanges converge absolutely on compact subsets of this
domain; $A_{u^2}(z)=O(u^{-2})$ at infinity. The constant terms cancel
by (KME5). The integral definition on the first line provides the
continuation of $G_\kappa$ to all of $\mathbb C\setminus[a,b]$.

Set
\begin{equation}\tag{KME11}
 U_\kappa(x)=V_\lambda(x)-2\int_a^b\log|x-t|\rho_\kappa(t)\,dt.
\end{equation}
It is continuous on $(0,\infty)$: the density is bounded and compactly
supported, and the integral of $|\log|x-t||$ over a small interval
about $x$ tends uniformly to zero with its length.  In the interior,
the derivative of its logarithmic potential is the principal value
integral. Subtracting $\rho_\kappa(x)$ on a symmetric neighbourhood of $x$
proves this directly, because the remaining numerator is $O(|x-t|)$.
Taking the real boundary value of (KME10) therefore gives $U_\kappa'=0$
there. The same subtraction proves that this real boundary value is
the indicated principal value. There is a finite constant $\ell_\kappa$ with
\begin{equation}\tag{KME12}
 U_\kappa=\ell_\kappa\ \hbox{on }[a,b],\qquad
 U_\kappa'(x)=\mathcal R(x)H_\kappa(x)
 \begin{cases}<0,&0<x<a,\\>0,&x>b.\end{cases}
\end{equation}
It follows that $U_\kappa(x)>\ell_\kappa$ at every positive point outside
$[a,b]$.  This proves the exterior inequality on the entire original
domain, including the interval adjacent to zero.

We give the energy positivity argument needed for global uniqueness.
For a finite signed real measure $\eta$ of total mass zero with
$|\log|x-y||$ integrable against $|\eta|\otimes|\eta|$, the identity
\begin{equation}\tag{KME13}
 -\log r=\frac12\int_0^\infty
           \frac{e^{-tr^2}-e^{-t}}t\,dt\quad(r>0)
\end{equation}
follows by differentiation in $r$ and evaluation at $r=1$.
Each truncated half-integral has absolute value at most $|\log r|$
and the same sign as $-\log r$. Dominated convergence and the zero
mass consequently give
\begin{align}\tag{KME14}
 -\iint\log|x-y|\,d\eta(x)d\eta(y)
 &=\frac12\int_0^\infty\frac1t
               \iint e^{-t(x-y)^2}\,d\eta(x)d\eta(y)\,dt\geq0,\\
 \iint e^{-t(x-y)^2}\,d\eta(x)d\eta(y)
 &=\frac1{\sqrt{4\pi t}}\int_{\mathbb R}
       e^{-\xi^2/(4t)}|\widehat\eta(\xi)|^2\,d\xi.\notag
\end{align}
The second identity follows by completing the square in the scalar
Gaussian Fourier integral and integrating it against the finite measure.
If $\eta\ne0$, its Fourier transform is nonzero somewhere and, being
continuous, makes the last integral strictly positive for every $t>0$.
For completeness, a measure with identically zero Fourier transform
vanishes: its convolution with every centred Gaussian is zero by that
same formula, and Gaussian approximate identities converge uniformly on
each compactly supported continuous test function.  Thus (KME14)
has equality only for $\eta=0$.

Let $\mu_\kappa(dx)=\rho_\kappa(x)dx$ and let $\mu$ have finite energy.
For $\eta=\mu-\mu_\kappa$ the absolute logarithmic integrability just used
holds. The self term for $\mu$ follows from its finite energy and
logarithmic moment; that for $\mu_\kappa$ follows from bounded compact
support. For the cross term, the negative logarithmic singularity
integrated against the bounded density is uniformly bounded, and its
positive part is bounded by $\log(1+x)+\log(1+y)$.
Expanding the two energies now gives
\begin{equation}\tag{KME15}
 \mathcal E_\lambda(\mu)-\mathcal E_\lambda(\mu_\kappa)
 =\int(U_\kappa-\ell_\kappa)\,d\mu
                -\iint\log|x-y|\,d\eta(x)d\eta(y)\geq0.
\end{equation}
The inequality is strict unless $\mu=\mu_\kappa$. Infinite energy
competitors cannot lower the finite energy of $\mu_\kappa$. Therefore the
explicit density above is the unique global minimizing probability
measure, and $F_\lambda=-\mathcal E_\lambda(\mu_\kappa)$ is its actual
energy value.

\subsection{Entropy and continuity on the actual parameter interval}
The endpoint proof gives smooth endpoints and $0<a(\kappa)<b(\kappa)<\infty$
on every compact $\kappa$ interval. On such an interval their distance from
zero and their separation have positive lower bounds. Uniform
differentiation in (KME6) proves that $H_\kappa(x)$ and its parameter
derivatives are continuous and bounded on all the resulting supports.
The positive integral in (KME6) also gives a positive lower bound there.
In the fixed coordinate $x=a+(b-a)y$, the probability measure is
\begin{equation}\tag{KME16}
 d\mu_\kappa(x)=\frac{(b-a)^2}{2\pi}\sqrt{y(1-y)}
           H_\kappa\bigl(a+(b-a)y\bigr)\,dy,\qquad 0<y<1.
\end{equation}
Its smooth coefficients are uniformly bounded on compact parameter
intervals, with each parameter derivative dominated by a constant times
$\sqrt{y(1-y)}$.  This proves continuity, and smoothness for smooth
observables, of all the integrals used below.

The probability reference in the original Gamma ensemble is
$d\omega(x)=\tfrac12x^{-1/2}e^{-\sqrt x}dx$.
The exact density ratio on the support is
\begin{equation}\tag{KME17}
 \frac{d\mu_\kappa}{d\omega}(x)=2\sqrt x\,e^{\sqrt x}\rho_\kappa(x).
\end{equation}
The factors other than the square-root endpoint factor are bounded
above and below by positive numbers uniformly for $\kappa$ in a compact
interval.  It follows from (KME16) that
\[
 \int\left|\log\frac{d\mu_\kappa}{d\omega}\right|\,d\mu_\kappa<\infty,
 \qquad
 \iint|\log|x-y||\,d\mu_\kappa(x)d\mu_\kappa(y)<\infty.
\]
Indeed the first absolute integral is bounded by a constant times
$\int_0^1\sqrt{y(1-y)}
(1+|\log y|+|\log(1-y)|)dy$, which is finite; for example
$\int_0^1y^{1/2}|\log y|dy=4/9$ by differentiation of
$\int_0^1y^pdy=1/(p+1)$. The second integral follows from bounded
compact support as above. These facts hold in particular uniformly
on the actual range $0\leq\lambda\leq1/4$.

\subsection{Explicit elementary endpoint integrals for the energy}
Define $d_{\rm end},m_{\rm end},D_s$ as in (KME4) and, for $z,s\geq0$, put
\begin{align}\tag{KME18}
 C_z&=\frac{m_{\rm end}+z+D_z}{2},\qquad
 \theta_s=\frac{d_{\rm end}}{m_{\rm end}+s+D_s},\\
 Q_s(z)&=(m_{\rm end}+s-D_s)\log C_z-d_{\rm end}\theta_z
                         -2D_s\log(1-\theta_z\theta_s).
 \notag
\end{align}
All logarithm arguments are positive: $0<\theta_s\leq\theta_0<1$.
The exact identity behind these formulas is
\begin{equation}\tag{KME19}
 Q_s(z)=\frac1\pi\int_a^b
          \frac{\sqrt{(b-x)(x-a)}}{x+s}\log(x+z)\,dx.
\end{equation}
To prove it, substitute $x=m_{\rm end}+d_{\rm end}\cos\phi$, and use the convergent series
\[
 \log(x+z)=\log C_z+
       2\sum_{j\geq1}\frac{(-1)^{j+1}\theta_z^j}{j}\cos(j\phi),
 \qquad
 \frac1{x+s}=\frac1{D_s}
       \left(1+2\sum_{j\geq1}(-\theta_s)^j\cos(j\phi)\right).
\]
Their averages are respectively $\log C_z$ and
$[\log C_z+2\log(1-\theta_z\theta_s)]/D_s$ after multiplication.
The average of $x\log(x+z)$ is $m_{\rm end}\log C_z+d_{\rm end}\theta_z$.
Finally divide
$(b-x)(x-a)/(x+s)=-x+2m_{\rm end}+s-D_s^2/(x+s)$ and combine these three
averages. This proves (KME19) with every constant in (KME18).

The logarithmic expectation, including the added mark parameter $z$, is
\begin{equation}\tag{KME20}
 M_\kappa(z):=\int\log(x+z)\,d\mu_\kappa(x)
       =\frac12\int_\kappa^\infty
                    \frac{Q_0(z)-Q_{u^2}(z)}{R_u}\,du.
\end{equation}
For proof, use (KME6) in the first expression, giving
$Q_0(z)/D_0-\tfrac12\int_\kappa^\infty Q_{u^2}(z)/R_u\,du$;
the first endpoint equation gives the stated form.
The exchanges are absolute, since $\log(x+z)$ is bounded on the
support, $Q_0(z)-Q_{u^2}(z)=O(u^2)$ near zero, and
$Q_{u^2}(z)=O(u^{-2})$ at infinity by (KME19).

For the logarithmic potential at the right endpoint define
\begin{align}\tag{KME21}
 C_{\rm edge}&=\frac{d_{\rm end}}2=\frac{b-a}{4},\\
 Q_s^{\rm edge}&=(m_{\rm end}+s-D_s)\log C_{\rm edge}
                            +d_{\rm end}-2D_s\log(1+\theta_s),\notag\\
 J_\kappa&=\frac12\int_\kappa^\infty
                  \frac{Q_0^{\rm edge}-Q_{u^2}^{\rm edge}}{R_u}\,du
       =\int\log(b-x)\,d\mu_\kappa(x).\notag
\end{align}
These identities follow by the same calculation with
$\log(b-x)=\log(d_{\rm end}/2)-2\sum_{j\geq1}\cos(j\phi)/j$.
One may justify this endpoint series by first replacing its coefficient
$1$ by $r<1$ and then taking $r\uparrow1$; the possible singularity is
bounded by a constant plus $|\log\phi|$, integrable on $[0,\pi]$.
The resulting averages are $\log(d_{\rm end}/2)$, $m_{\rm end}\log(d_{\rm end}/2)-d_{\rm end}$, and
$[\log(d_{\rm end}/2)+2\log(1+\theta_s)]/D_s$, respectively.
This proves the displayed $Q_s^{\rm edge}$ and then $J_\kappa$ by the
same endpoint cancellation as (KME20). All the final integrals converge
absolutely, including the logarithmic singularity at $x=b$.

There is now an explicit energy formula involving only the uniquely
specified endpoints and real one-dimensional integrals:
\begin{equation}\tag{KME22a}
 \boxed{F_\lambda=J_\kappa-\frac12V_\lambda(b)-3-\kappa
                       +M_\kappa(0)+\frac \kappa2M_\kappa(\kappa^2).}
\end{equation}
Here is a direct derivation of the exact constant $3+\kappa$.
Push the minimizing probability measure forward by the explicit
coordinate dilation $x\mapsto r x$ for $r>0$. Differentiating its energy
at its minimizing value $r=1$ is legitimate on its compact positive
support and gives $\int xV'_\lambda(x)d\mu_\kappa(x)=1$: the logarithmic
double integral contributes exactly $-\log r$ to the energy.
By (KME2) this is
$\int\sqrt x\arctan(\sqrt x/\kappa)d\mu_\kappa(x)=3$.
Taking the expectation of the first formula in (KME2) now gives
$\int V_\lambda d\mu_\kappa=6+2\kappa-2M_\kappa(0)-\kappa M_\kappa(\kappa^2)$.
Finally $\ell_\kappa=V_\lambda(b)-2J_\kappa$, and the integrated Euler equality
gives $F_\lambda=-\tfrac12(\int V_\lambda d\mu_\kappa+\ell_\kappa)$.
These identities prove (KME22a).

An equivalent representation, useful for checking all the additive
constants directly against the Stieltjes integral, is
\begin{equation}\tag{KME22}
 \boxed{F_\lambda=J_\kappa+M_\kappa(0)-\frac12V_\lambda(b)
      -\frac12\int_\kappa^\infty\bigl(M_\kappa(u^2)-2\log u\bigr)du
                          +\kappa\log \kappa-\kappa.}
\end{equation}
At $\kappa=0$ set $\kappa\log \kappa=0$.
To prove the formula without an unspecified energy constant, first
observe the scalar identity
\begin{equation}\tag{KME23}
 V_\lambda(x)=\int_\kappa^\infty\log(1+x/u^2)\,du
                           -2\log x-2\kappa\log \kappa+2\kappa.
\end{equation}
The identity $\int_0^\infty\log(1+x/u^2)du=\pi\sqrt x$
follows by differentiation in $x>0$ and then its zero limit as
$x\downarrow0$ (or by $u=\sqrt x v$ and one integration by parts).
Subtracting the integral over $[0,\kappa]$ proves (KME23), including the
integral $\int_0^\kappa2\log u\,du=2\kappa\log \kappa-2\kappa$.
Its expectation is absolutely finite: near zero $|\log u|$ is
integrable, and at infinity
$0\leq M_\kappa(u^2)-2\log u
=\int\log(1+x/u^2)d\mu_\kappa\leq b/u^2$.
Since $\ell_\kappa=V_\lambda(b)-2J_\kappa$ by (KME11)--(KME12), integration
of the Euler equality gives
$F_\lambda=-\tfrac12(\int V_\lambda d\mu_\kappa+\ell_\kappa)$.
Substituting (KME23) proves exactly (KME22).

The difference of the actual energies also has the following exact form:
\begin{equation}\tag{KME24}
 \boxed{\frac{dF_\lambda}{d\lambda}=2M_{2\lambda}(4\lambda^2),
 \qquad F_\lambda-F_0=2\int_0^\lambda M_{2t}(4t^2)\,dt.}
\end{equation}
Here the subscript of $M_\kappa$ is the lower Stieltjes endpoint $\kappa$, so
the parameter relation in the formula is explicit.
For a direct proof, the difference
$V_{\lambda+h}(x)-V_\lambda(x)
=-2\int_\lambda^{\lambda+h}\log(x+4t^2)dt$
is exact. Testing the two unique minimizers in each other's variational
problem bounds the difference quotient of $F$ between the integrals
of $2h^{-1}\int_\lambda^{\lambda+h}\log(x+4t^2)dt$ against these
two measures. Formula (KME16) and compact positive support bounds show
that both bounds tend to $2M_{2\lambda}(4\lambda^2)$.
This includes the right derivative at zero. Continuity of this
derivative and the fundamental theorem of calculus prove the second
equation in (KME24).

The formulas above apply, in particular, to the actual macroscopic
choice $K=q+1$, $L=q/4$ when $q$ is a positive multiple of four.
Its limiting mark parameter is exactly $\lambda=1/4$ and the lower
Stieltjes endpoint is exactly $\kappa=1/2$. The energy in (KME22) or (KME24)
is computed for the entire changed potential (KME1), with its additive
constants and original $-2\log x$ term preserved.

Finally the actual energy is convex as a function of $\lambda$.
For each fixed compactly supported probability measure of finite
logarithmic energy on $(0,\infty)$,
its objective in (KME3) has second parameter derivative
$16\lambda\int(x+4\lambda^2)^{-1}d\mu(x)\geq0$.
The supremum of these convex functions is $F_\lambda$: its proved
minimizer belongs to this class for every parameter. Therefore the
defining convexity inequality passes to that supremum, and (KME24)
gives the unconditional tangent estimate
\begin{equation}\tag{KME25}
 F_\lambda-F_0\geq2\lambda M_0(0)\qquad(\lambda\geq0).
\end{equation}
All the quantities on the right are given by (KME18)--(KME20) at the
uniquely constructed zero-parameter endpoints.
